\documentclass[a4paper,11pt]{article}
\usepackage{amsmath, amssymb}

\begin{document}

\title{Fourier Transform Solution for RLC Circuit}
\author{}
\date{}
\maketitle

\section*{Introduction}
This document provides the Fourier transform solution for the RLC circuit with a driven AC current. We derive the governing equations and analyze the circuit in the frequency domain using the Fourier transform technique.

\section*{Circuit Description}
Consider a series RLC circuit driven by an AC voltage source:
\[
V_{\text{in}}(t) = V_0 e^{i\omega t}.
\]
The components of the circuit are:
\begin{itemize}
    \item Resistor: \(R\),
    \item Inductor: \(L\),
    \item Capacitor: \(C\).
\end{itemize}

The current through the circuit is denoted by \(I(t)\).

\section*{Governing Equation}
Applying Kirchhoff's Voltage Law (KVL):
\[
L \frac{d^2q}{dt^2} + R \frac{dq}{dt} + \frac{q}{C} = V_{\text{in}}(t).
\]
Here, \(q(t)\) is the charge on the capacitor and \(I(t) = \frac{dq}{dt}\) is the current. Substituting \(I(t)\):
\[
L \frac{dI}{dt} + R I + \frac{1}{C} \int I \, dt = V_{\text{in}}(t).
\]
Taking the Fourier transform of both sides:
\[
\mathcal{F}\left[L \frac{dI}{dt} + R I + \frac{1}{C} \int I \, dt \right] = \mathcal{F}\left[V_{\text{in}}(t)\right].
\]
Using the properties of the Fourier transform:
\begin{itemize}
    \item \(\mathcal{F}\left[\frac{dI}{dt}\right] = i\omega \hat{I}(\omega)\),
    \item \(\mathcal{F}\left[\int I \, dt\right] = \frac{\hat{I}(\omega)}{i\omega}\),
    \item \(\mathcal{F}\left[I(t)\right] = \hat{I}(\omega)\),
\end{itemize}
where \(\hat{I}(\omega)\) is the Fourier transform of \(I(t)\).

The equation in the frequency domain becomes:
\[
i\left(L i\omega \hat{I}(\omega) + R \hat{I}(\omega) + \frac{1}{C} \frac{\hat{I}(\omega)}{i\omega}\right) = \hat{V}_{\text{in}}(\omega).
\]

Simplify:
\[
\hat{I}(\omega) \left(R + i\omega L - \frac{i}{\omega C}\right) = \hat{V}_{\text{in}}(\omega).
\]

\section*{Solution in Frequency Domain}
The impedance \(Z(\omega)\) of the circuit is given by:
\[
Z(\omega) = R + i\omega L - \frac{i}{\omega C}.
\]
Thus, the Fourier transform of the current is:
\[
\hat{I}(\omega) = \frac{\hat{V}_{\text{in}}(\omega)}{Z(\omega)}.
\]
For a sinusoidal input voltage \(V_{\text{in}}(t) = V_0 e^{i\omega t}\), we have:
\[
\hat{V}_{\text{in}}(\omega) = V_0 \delta(\omega - \omega_0).
\]
Substituting:
\[
\hat{I}(\omega) = \frac{V_0 \delta(\omega - \omega_0)}{R + i\omega L - \frac{i}{\omega C}}.
\]

\section*{Solution in Time Domain}
The current in the time domain is obtained by taking the inverse Fourier transform:
\[
I(t) = \mathcal{F}^{-1}\left[\hat{I}(\omega)\right].
\]
Substitute \(\hat{I}(\omega)\):
\[
I(t) = \mathcal{F}^{-1}\left[\frac{V_0 \delta(\omega - \omega_0)}{R + i\omega L - \frac{i}{\omega C}}\right].
\]
Since \(\delta(\omega - \omega_0)\) selects the frequency \(\omega_0\):
\[
I(t) = \frac{V_0}{Z(\omega_0)} e^{i\omega_0 t}.
\]

The magnitude of the current is:
\[
|I(t)| = \frac{V_0}{\sqrt{R^2 + \left(\omega_0 L - \frac{1}{\omega_0 C}\right)^2}}.
\]
The phase shift \(\phi\) is:
\[
\phi = \tan^{-1}\left(\frac{\omega_0 L - \frac{1}{\omega_0 C}}{R}\right).
\]
Thus, the time-domain current is:
\[
I(t) = \frac{V_0}{\sqrt{R^2 + \left(\omega_0 L - \frac{1}{\omega_0 C}\right)^2}} \cos(\omega_0 t + \phi).
\]

\section*{Conclusion}
Using the Fourier transform, we derived the current \(I(t)\) for a driven RLC circuit. This approach simplifies the analysis by handling the problem algebraically in the frequency domain, and the inverse Fourier transform retrieves the time-domain solution.

\end{document}
